\section{Once-counted distinct-port three-block cancellation}
\label{sec:three-block-unmarking}
For the unchanged frozen propagator, put $P_{xy}=|U_{xy}|^2$ and
$J_{xy}=N^{-1}$. The position-chain theorem supplies, for all sufficiently
large odd $N$, $\max P_{xy}\leq C/N$ and
$\|P^L-J\|\leq 2^{-L}$ for every integer $L\geq0$.
Let $T_m$ be the set of rooted ordered coordinate-word pairs obtained by
cutting the first cyclic word into three nonempty blocks $A,B,C$ with
pairwise distinct initial vertex values, then taking a cyclic rotation of
$ACB$ for the second word. This is a set: each ordered pair is counted once,
even if several cuts produce it. Write
$\mathcal A(x)=\prod_{j\bmod m}U_{x_jx_{j+1}}$.

\paragraph{Theorem (BGS-TRACE-015).}
For every smooth $w$ compactly supported in $(0,\infty)$,
\[
 \frac1{N^2}\sum_m w(m/N)
 \sum_{(x,y)\in T_m}\mathcal A(x)\overline{\mathcal A(y)}\longrightarrow0.
\]
For $K=\lceil\log N\rceil$ and support in $[a,b]\subset(0,\infty)$,
the error is bounded by
\[
 O_w(K^5/N)+O_w\!\left(N^8(C/N)^{K/2}\right)
 +O_w\!\left(N^5 2^{-(aN-3)/3}\right).
\]
The startup threshold inherited from the position-chain theorem is not
made effective.

\paragraph{Proof mechanism.}
In the marked three-link contraction, replacing any link by $J$ gives
zero by column orthogonality and distinct ports. Using the longest link
therefore suppresses the full marked sum exponentially, including zero
internal link lengths. Passing to a once-counted family is a separate
combinatorial step. Words with repeated length-$K$ substrings impose
at least $K/2$ equality constraints and have superpolynomially small
exceptional absolute weight. For the remaining words, two blocks of
length at least $4K$ force a unique cut triple, using uniquely located
common substring anchors. Ambiguity is consequently confined to two
short blocks. Competing cuts lie in an $O(K)$ neighborhood, giving
$O(K^3)$ alternatives and $O(K^2)$ original length choices.
A rank-one relative positional permutation would be a transposition;
conjugating predecessor cycles then forces equal-valued positions to be
two cut heads, contradicting distinct ports. Thus every competing
representation imposes at least two equalities. Exposing the three starts
first pays each equality through a transition bounded by $C/N$. After the
two trace roots and band normalization, the exceptional mass is
$O(K^5/N)$. Every nonexceptional pair has exactly three marked
representations, corresponding to cyclic block labels. The canonical
proof gives all equality, cyclic-root and overlap counts explicitly.

\paragraph{Boundary.}
This is a deterministic, once-counted, fixed-complexity family estimate.
The combinatorial theorem uses only unitarity, flat entries and the stated
position-chain mixing bound; actual-model input enters through BGS-FM-021.
It is compatible with spectrally nonuniversal families satisfying those
matrix estimates.
It does not cover repeated ports, arbitrary reconnections, growing
encounter complexity or higher joint traces. The complete signed remainder,
including restoration of the super-Heisenberg plateau, remains open;
no BGS or local-statistics theorem follows from this family alone.
